\scanpage{91}

\deleted{\unsure{La formule}} Enfin, si $\xi=(\xi_i)$ est une suite quelconque de nbres réels \deleted{positifs}, posons $x_n=(\xi_1,\ldots\xi_n,0,\ldots0)$, alors $N(x_n)$ \deleted{a une limite \unsure{quand} \illegible{1} de $n$}, \inserted{croît avec $n$, et} on pose $N(\xi)=N((\xi_i))=\lim_n N(x_n)$, nb fini ou infini $\geq0$.

\deleted{\unsure{Ainsi|Alors}} La formule (25) est un cas particulier du théorème, car $N$ \deleted{\unsure{jouit des propriétés}}\inserted{satisfait aux} conditions envisagées pour $f$ dans cet énoncé. La formule (26) en résulte par un passage à la limite.

\uline{Th 3} \uline{Deuxième théorème de convexité dans les espaces de Hilbert}

Soi\deleted{\unsure{ent}}\unsure{ent} $E,F,G$ trois espaces de Hilbert, soit $u$ une application linéaire \deleted{continue}\inserted{compacte} de $E$ dans $F$, \deleted{\unsure{$E$|et}} $v$ une application linéaire \deleted{continue}\inserted{compacte} de $F$ dans $G$, \deleted{soit $u$ un entier positif} enfin soit $f(s_1,\ldots s_n)$ une fonction continue \inserted{symétrique et croissante} de $n$ arguments $s_i\geq0$, telle que \deleted{la fonction} $f(e^{t_1},\ldots,e^{t_n})$ soit fonction convexe de $(t_i)\in R^n$. Alors on a

\[
f(\rho_1(\unsure{uv|vu}),\ldots\rho_n(uv))\leq f(\rho_1(u)\rho_1(v),\ldots,\rho_n(u)\rho_n(v))\tag*{(29)}
\]

En effet, il suffit de \deleted{conjuguer}\inserted{\unsure{l'appliquer|appliquer}} le \deleted{lemme}\inserted{th. 1} \deleted{fondamental, \unsure{corollaire} \illegible{1} \unsure{condition}} aux deux suites $(\rho_i(uv))_{i\leq n}$ et $(\rho_i(u)\rho_i(v))_{\unsure{i\leq n|i=n}}$. Les conditions \deleted{du corollaire 1 de ce lemme}\inserted{de l'hyp. b, n° 2°) du th 1} sont satisfaites, \deleted{\unsure{en ef}} en vertu du N° \unsure{?|3}, prop. \unsure{2|3}, corollaire 3. \deleted{\unsure{(et} \illegible{2}\unsure{)}.}

\uline{Corollaire 1} Soient $u$ et $v$ \deleted{comme} dans le th. 3, soit $f(r)$ une fonction continue croissante de $r\geq0$ telle que $f(e^t)$ soit fonction convexe de $t\in R$. Alors, pour \deleted{tout \unsure{nbre} entier} $n>0$, on a

\[
\sum_{i=1}^n f(\rho_i(uv))\leq\sum_{i=1}^n f(\rho_i(u)\rho_i(v))\tag*{(30)}
\]

\deleted{d'où} d'où

\[
\sum_{\unsure{i|j}} f(\rho_i(uv))\leq\sum_i f(\rho_i(u)\rho_i(v))\tag*{(31)}
\]

(les sommes étant étendues à tous les indices $i$)

\deleted{\unsure{$l$} \illegible{1} $\unsure{(\rho_i(uv))|(\rho_i(vu))}$ associée \illegible{2} \unsure{des suites} $(\rho_i(u))$}

\deleted{$(\rho_i(v))$ associée à \unsure{l'opér.} $v$, la $=$ \unsure{majoration|majorations} \deleted{\unsure{pour tout}} que celle qui vaudrait si on avait $\rho_i(uv)=\rho_i(u)\rho_i(v)$, donc si $u$ et $v$ étaient des opérateurs de multiplication dans $l^{\unsure{2|?}}$ par les \unsure{suites} $(\rho_i(u))$ resp.t $(\rho_i(v))$. \deleted{\unsure{En particularisant}} \inserted{Notons qu'on pourrait aussi} \deleted{\unsure{ce dernier, il}} donner un corollaire \deleted{qui} analogue \deleted{\unsure{Théorème 4} \unsure{Soient}} au corollaire 4 du th. 1, mais en laissons l'énoncé au lecteur.}

\deleted{Dans l'inégalité (39), \unsure{prenons|on peut faire} $p$ fixé et \unsure{très} \illegible{1} $q$ vers $0$, \unsure{$u$} donné par (30), \illegible{1} et on obtient à la limite}

% archive sheet number: 92
